\section{Introduction}

\subsection{Contexte et problématique}

Dans le cadre du module de Calcul Embarqué du semestre 6 (ING3) à l'ECE Paris, il est demandé de concevoir un système de reconnaissance vocale fondé sur une intelligence artificielle s'exécutant \emph{intégralement} sur une carte Arduino Due (SAM3X8E, \SI{84}{\mega\hertz}, \SI{96}{\kibi\byte} de SRAM). Le projet, nommé \textbf{NeuralSpeech}, s'inscrit dans la continuité pédagogique des modules de traitement numérique du signal, d'architecture microcontrôleur et d'apprentissage automatique.

La ligne rouge à ne pas franchir est formelle : \emph{la reconnaissance doit se faire sur l'Arduino Due impérativement}. Un projet fonctionnel sur PC mais incapable de tourner sur le microcontrôleur entraîne un malus forfaitaire de -12 points sur la note finale.

\subsection{Objectifs du projet}

L'objectif est de réaliser une chaîne complète de traitement audio, depuis la captation analogique jusqu'à la classification temps réel :
\begin{enumerate}
    \item Numériser un signal audio à \SI{32}{\kilo\hertz} via l'ADC intégré.
    \item Filtrer numériquement puis sous-échantillonner à \SI{8}{\kilo\hertz}.
    \item Valider l'enregistrement par export sur Audacity.
    \item Extraire les coefficients MFCC caractéristiques du timbre vocal.
    \item Entraîner un réseau de neurones convolutionnel en Python.
    \item Embarquer l'inférence sur l'Arduino Due et signaler le mot détecté par des LEDs.
\end{enumerate}

\subsection{Cadre et contraintes (ET1 à ET9)}

Le sujet impose neuf exigences techniques (ET) non négociables, récapitulées au tableau~\ref{tab:et}.

\begin{table}[H]
\centering
\small
\begin{tabularx}{\linewidth}{@{}l X@{}}
\toprule
\textbf{Exigence} & \textbf{Libellé} \\
\midrule
\etref{1} & ADC \SI{32}{\kilo\hertz} via interruption timer, échantillonnage fixe et précis \\
\etref{2} & Atténuation $\geq \SI{30}{\decibel}$ au-dessus de \SI{4}{\kilo\hertz} \\
\etref{3} & Filtrage $< \SI{31}{\micro\second}$ par échantillon \\
\etref{4} & Enregistrement du mot « Électronique » restituable sur Audacity \\
\etref{5} & Frames de 256 échantillons avec recouvrement (hop \texttt{128}) \\
\etref{6} & 13 MFCCs générés à partir d'un enregistrement d'\SI{1}{\second} \\
\etref{7} & MSE $< 0{,}05$ sur données de test après entraînement \\
\etref{8} & Dataset $\geq$ 100 éléments (50 par mot, $\geq 2$ mots) \\
\etref{9} & $< 5\%$ d'erreur sur $\geq 10$ mots prononcés, robustesse au bruit de fond \\
\bottomrule
\end{tabularx}
\caption{Exigences techniques imposées par le cahier des charges.}
\label{tab:et}
\end{table}

\subsection{Matériel}

\begin{itemize}
    \item \textbf{Arduino Due} (microcontrôleur Atmel SAM3X8E, ARM Cortex-M3, \SI{84}{\mega\hertz}, \SI{96}{\kibi\byte} SRAM, \SI{512}{\kibi\byte} Flash, ADC 12 bits, deux DAC 12 bits)
    \item \textbf{Microphone MAX9814} (amplificateur electret avec contrôle automatique de gain)
    \item Bouton poussoir, LEDs de signalisation, breadboard, résistances
    \item Oscilloscope Siglent SDS 1102CML+ et générateur de fonctions Siglent SDG 805 pour la validation
\end{itemize}

\subsection{Organisation du rapport}

Ce rapport suit l'architecture fonctionnelle du système. Après une présentation globale de l'architecture (section~\ref{sec:architecture}), chaque fonction principale (FP1 à FP6) fait l'objet d'une section dédiée organisée en trois volets : conception, implémentation et validation. La section~\ref{sec:gestion-projet} détaille la gestion de projet (répartition des tâches, difficultés rencontrées), et la section~\ref{sec:conclusion} dresse le bilan et les perspectives d'évolution.

\newpage
